\SourceLeaf{203}\Needspace{10\baselineskip}\begin{center}
[4.]
\end{center}

[Bei \textsc{Lübsen}, \emph{Ausf\"uhrliches Lehrbuch der Elementar-Geometrie}, Hamburg 1851, Seite 11 heißt es:]

\{\textls{Erklärung:} Eine gerade Linie ist diejenige, welche nicht aus ihrer Lage kommt, indem sie sich um ihre beiden festen Endpunkte dreht\HistoricalNote{*)}{\{*)}{So ungef\"ahr h\"orten wir einmal \textsc{Gauss} bei der Erkl\"arung des Fernrohrs und dessen richtigen Gebrauchs den Begriff der geraden Linie festsetzen. Diese Erkl\"arung ist theoretisch fruchtbar, wie die gleich daraus folgenden S\"atze zeigen; ausserdem ist das angegebene Merkmal praktisch wichtig, z.\,B. bei der Justirung eines Fernrohrs, richtigen Bohrung eines Cylinders etc.\}}.\}



\begin{center}
[5.]

\textnormal{THEORIE DES VORTRAGS VON LEHREN, DIE RAUMVERHÄLTNISSE BETREFFEN.}
\end{center}

1) Alle Gegenstände, die bei einem Satz von Belang sind, \textls{zeichne} man und bezeichne sie durch Buchstaben oder Ziffern, Punkte durch einen, Linien durch 2, Flächen durch 3 und Körper gegebenenfalls durch 4, wobei man auch hier oft mit einem auskommt.

2) Räume, die sich relativ zueinander bewegen können und deren Bewegung bei der Aufgabe ohne Stetigkeit in Frage kommt, zeichne man gesondert.

3) Auch die relevanten Größenbeziehungen bezeichne man durch Buchstaben.

4) Man stelle das, was man beweisen will, als Gleichung dar und gehe schrittweise auf ihre Gründe zurück.

So wird die ganze Aufgabe gleichsam in ihre Elemente zerlegt, die man danach in umgekehrter Reihenfolge wieder zusammensetzt. Dabei erkennt man leicht, was sich im eigentlichen Vortrag zusammenfassen lässt.

